% ---------------------------------------------------------------------
% Section 2: Monotone POMDPs
% Source: notes/source-material/LRMQG.tex, Sections "Setup", "Beliefs",
% "The Dynamic Program"; guiding text drafted by Claude.
% Propositions stated without proof per the authors' instruction;
% surrounding text points the reader to the relevant literature.
% ---------------------------------------------------------------------
\section{Monotone POMDPs}\label{sec:model}

\subsection{States, Actions, and Beliefs}\label{sec:setup}

We consider a discrete-time, finite-horizon POMDP. The state of the
world is denoted $\theta \in \Theta$, where $\Theta \subseteq \Re$ is a
fully ordered set; for example, $\theta$ may represent the benefit of a
new technology in a technology adoption model or the condition of a
machine in a machine replacement model. In each period, the decision
maker (DM) chooses an action $a$ from a finite and fully ordered set
$A = \{a_0, a_1, \ldots, a_N\}$, where $a_0 \leq a_1 \leq \cdots \leq
a_N$. The DM cannot observe the state directly. Her beliefs about
$\theta$ are described by a probability distribution; for ease of
notation, we assume that this distribution is continuous with a density
$\pi$ over $\Theta$. For discrete state spaces, we can interpret $\pi$
as a probability mass function and replace integrals with sums
\citep{ulu2009uncertainty}.

We index periods by the number of periods remaining, $k$. The state may
change over time: if the state in period $k$ is $\theta_k$ and the DM
takes action $a$, the next-period state $\theta_{k-1}$ is drawn from the
conditional distribution $g(\theta_{k-1}|\theta_k, a)$. Given a prior
$\pi$ on the period-$k$ state, the DM's \emph{next-period prior} is
\begin{equation}
\eta(\theta_{k-1}; \pi, a)
  = \int_{\theta_k \in \Theta} g(\theta_{k-1}|\theta_k, a)\,
    \pi(\theta_k) \dif \theta_k \ . \label{eq:eta}
\end{equation}
Information arrives through a \emph{signal process} $\iX$ that generates
signals $x \in X$, where $X$ is fully ordered, with likelihood function
$L_\iX(x|\theta, a)$ describing the probability of observing signal $x$
given (next-period) state $\theta$ and action $a$. The \emph{predictive
distribution} for signals and the \emph{posterior distribution} on the
state given a signal are then
\begin{align}
f_\iX(x;\pi, a) &= \int_\theta L_\iX(x|\theta,a)\, \eta(\theta; \pi, a)
  \dif \theta \ , \label{eq:predictive}\\
\Pi_\iX(\theta; \pi, x, a) &=
  \frac{L_\iX(x|\theta,a)\, \eta(\theta; \pi, a)}{f_\iX(x;\pi, a)} \ .
  \label{eq:posterior}
\end{align}
Thus, within a period, beliefs evolve from the prior $\pi$ to the
next-period prior $\eta(\,\cdot\,; \pi, a)$ through the state
transition, and then to the posterior $\Pi_\iX(\,\cdot\,; \pi, x, a)$
after the signal $x$ is observed; this posterior serves as the prior for
the next period. Because the distributions themselves will serve as
arguments of value functions, we will frequently suppress the state
argument and write, e.g., $\eta(\pi,a)$, $f_\iX(\pi,a)$, and
$\Pi_\iX(\pi,x,a)$ when we consider these distributions as functions of
the prior $\pi$, the signal $x$, and the action $a$.

When studying monotonicity properties involving distributions, we focus
on likelihood-ratio (LR) dominance as a stochastic order.

\begin{definition}[LR order and MLR property]\label{def:LR}\quad
\begin{enumerate}
\item A distribution $\pi_2$ \emph{LR-dominates} $\pi_1$
  ($\pi_2 \LR \pi_1$) if for all $\theta_2 \geq \theta_1$,
  \begin{equation}
  \frac{\pi_2(\theta_2)}{\pi_1(\theta_2)} \geq
  \frac{\pi_2(\theta_1)}{\pi_1(\theta_1)} \ . \label{eq:LRdef}
  \end{equation}
\item $L_\iX(x|\theta, a)$ satisfies the \emph{MLR property} if
  $L_\iX(\,\cdot\,|\theta_2, a) \LR L_\iX(\,\cdot\,|\theta_1, a)$
  whenever $\theta_2 \geq \theta_1$. Similarly, $g(\theta_{k-1}|
  \theta_k, a)$ satisfies the MLR property if
  $g(\,\cdot\,|\theta_k^2, a) \LR g(\,\cdot\,|\theta_k^1, a)$ whenever
  $\theta_k^2 \geq \theta_k^1$.
\end{enumerate}
\end{definition}

The ratios in \eqref{eq:LRdef} may be interpreted as being infinite when
the denominators are zero; equivalently, $\pi_2 \LR \pi_1$ if
$\pi_2(\theta_2)\pi_1(\theta_1) \geq \pi_2(\theta_1)\pi_1(\theta_2)$
for all $\theta_2 \geq \theta_1$. LR dominance implies first-order
stochastic dominance (FOSD) and, unlike FOSD, LR dominance survives
Bayesian updating and the state transition process, provided the
likelihoods and state transitions satisfy the MLR property; see
\citet{ulu2009uncertainty} for a discussion of this distinction and
examples where FOSD fails to be preserved. The following proposition
records these properties in our setting. The results are well known
(see, e.g., \citealt{whitt1979note}; \citealt{lovejoy1987some},
Lemma~1.2, for the finite-state case; and \citealt{ulu2009uncertainty},
Propositions~3.2 and 3.4 and \S8), and we state them without proof.

\begin{proposition}[Monotonicity of beliefs]\label{prop:monotone_beliefs}
For any action $a$, if $L_\iX(x|\theta,a)$ and
$g(\theta_{k-1}|\theta_k,a)$ satisfy the MLR property, then:
\begin{enumerate}
\item $\pi_2 \LR \pi_1$ implies $\eta(\pi_2,a) \LR \eta(\pi_1,a)$,
  $f_\iX(\pi_2,a) \LR f_\iX(\pi_1,a)$, and
  $\Pi_\iX(\pi_2, x, a) \LR \Pi_\iX(\pi_1, x, a)$ for all $x$.
\item For any $\pi$, $x_2 \geq x_1$ if and only if
  $\Pi_\iX(\pi,x_2,a) \LR \Pi_\iX(\pi,x_1,a)$.
\end{enumerate}
\end{proposition}

Part (i) says that more optimistic priors lead to more optimistic
next-period priors, signal distributions, and posteriors; part (ii) says
that more favorable signals lead to more optimistic posteriors. This
``survival'' of the LR order through Bayesian updating and state
transitions is the key property that drives the recursive arguments in
the dynamic program below.
\CU{Mention more properties of the order? e.g., single-crossing
preservation/dominance??}

\subsection{The Dynamic Program}\label{sec:dp}

The DM earns a reward $r(\theta, a)$ when taking action $a$ in state
$\theta$; given beliefs $\pi$, the expected reward is
\begin{equation}
r(\pi,a) = \int_{\theta \in \Theta} r(\theta,a)\, \pi(\theta)
  \dif \theta \ . \label{eq:expreward}
\end{equation}
Rewards are discounted with a discount factor $\delta \in [0,1]$. The
DM's optimal value function with $k$ periods remaining is defined by the
dynamic programming recursion: let $V_{\iX,0}^*(\pi) = 0$ and, for
$k > 0$,
\begin{align}
V_{\iX,k}(\pi,a) &= r(\pi,a) +
  \delta \ev{V_{\iX,k-1}^*(\Pi_\iX(\pi,\tilde{x},a))} \ ,
  \label{eq:Vka}\\
V_{\iX,k}^*(\pi) &= \max_{a\in A} \left\{ V_{\iX,k}(\pi,a) \right\} \ ,
  \label{eq:Vk}
\end{align}
where the expected continuation value is taken with respect to the
predictive distribution,
\begin{equation}
\ev{V_{\iX,k-1}^*(\Pi_\iX(\pi,\tilde{x},a))}
  = \int_{x\in X} V_{\iX,k-1}^*(\Pi_\iX(\pi,x,a))\,
    f_\iX(x; \pi, a) \dif x \ . \label{eq:contval}
\end{equation}
An \emph{action selection} $\alpha_{k}(\pi)$ chooses an action
$a \in A$ in period $k$ for a given prior $\pi$; a \emph{policy}
$\boldsymbol{\alpha} = (\alpha_{k}, \ldots, \alpha_{0})$ is a sequence
of action selections. An action selection $\alpha_{k}^*$ is optimal if
$V_{\iX,k}^*(\pi) = V_{\iX,k}(\pi,\alpha_{k}^*(\pi))$ for all $\pi$,
and a policy $\boldsymbol{\alpha}^*$ is optimal if each of its action
selections is optimal.

We say a function $u(\pi)$ defined on distributions is
\emph{LR-increasing} if $u(\pi_2) \geq u(\pi_1)$ whenever
$\pi_2 \LR \pi_1$. Our first structural result states that, with
increasing rewards and MLR likelihoods and transitions, the value
functions are LR-increasing: LR improvements in the DM's beliefs make
the DM better off. The proof is by induction on $k$, mirroring the
proof of Proposition~1 of \citet{lovejoy1987some} in the finite-state
case and Proposition~3.6 of \citet{ulu2009uncertainty} in the technology
adoption model: expected rewards \eqref{eq:expreward} are LR-increasing
because $r(\theta,a)$ is increasing in $\theta$ and LR dominance implies
FOSD; the expected continuation value \eqref{eq:contval} is
LR-increasing by Proposition~\ref{prop:monotone_beliefs} and the
induction hypothesis; and the maximum \eqref{eq:Vk} of LR-increasing
functions is LR-increasing. We omit the details.

\begin{proposition}[Monotonicity of value functions]\label{prop:monotone_values}
If, for all actions $a$, $r(\theta, a)$ is increasing in $\theta$ and
$L_\iX(x|\theta,a)$ and $g(\theta_{k-1}|\theta_k,a)$ satisfy the MLR
property, then:
\begin{enumerate}
\item $r(\pi,a)$ and $V_{\iX,k}(\pi,a)$ are LR-increasing for all $a$
  and $k$, and
\item $V_{\iX,k}^*(\pi)$ is LR-increasing for all $k$.
\end{enumerate}
\end{proposition}

\begin{remark}[Comparison with \citealt{lovejoy1987some}]\label{rem:lovejoy}
Proposition~\ref{prop:monotone_values} is a general-state-space analog
of Proposition~1 of \citet{lovejoy1987some}, and the assumptions map
to his as follows. We assume that the state transition process
$g(\theta_{k-1}|\theta_k,a)$ satisfies the MLR property; Lovejoy's
Lemma~1.3(1) shows that this (his TP$_2$ condition on the transition
matrix) is a sufficient condition for the monotonicity of next-period
priors required in his Lemma~1.2(2)---a condition that is, in his
finite-state setting, also necessary. Lovejoy's Lemma~1.3(2) provides
conditions under which next-period priors are LR-increasing in the
action; neither Lovejoy's Proposition~1 nor our
Proposition~\ref{prop:monotone_values} requires this assumption.
\end{remark}

\subsection{Monotone Optimal Policies}\label{sec:policies}

To establish the monotonicity of optimal policies, we introduce the
notion of increasing differences, also known as supermodularity.

\begin{definition}[Increasing differences]\label{def:incdiff}\quad
\begin{enumerate}
\item A function $u(\theta,a)$ has \emph{increasing differences} if
  $u(\theta_2,a_2)-u(\theta_2,a_1) \geq u(\theta_1,a_2)-u(\theta_1,a_1)$
  whenever $\theta_2 \geq \theta_1$ and $a_2 \geq a_1$.
\item A function $u(\pi,a)$ has \emph{LR-increasing differences} if
  $u(\pi_2,a_2)-u(\pi_2,a_1) \geq u(\pi_1,a_2)-u(\pi_1,a_1)$ whenever
  $\pi_2 \LR \pi_1$ and $a_2 \geq a_1$.
\item The state transitions defined by $L_\iX(x|\theta,a)$ and
  $g(\theta_{k-1}|\theta_k,a)$ satisfy a \emph{stochastic version of
  increasing differences} if
  \begin{equation}
  \ev{u(\Pi_\iX(\pi_2,\tilde{x},a_2))}-\ev{u(\Pi_\iX(\pi_2,\tilde{x},a_1))}
  \geq
  \ev{u(\Pi_\iX(\pi_1,\tilde{x},a_2))}-\ev{u(\Pi_\iX(\pi_1,\tilde{x},a_1))}
  \label{eq:stochincdiff}
  \end{equation}
  for all LR-increasing functions $u(\pi)$, whenever $\pi_2 \LR \pi_1$
  and $a_2 \geq a_1$.
\end{enumerate}
\end{definition}

Note that the functions $u(\pi)$ in part (iii) are assumed to be
LR-increasing; this could be replaced by other conditions, e.g.,
LR-increasing and convex, or many other choices.

The next lemma records two standard consequences of increasing
differences. Part (i) shows that increasing differences in the state
carry over to LR-increasing differences in beliefs when taking
expectations. Part (ii) is the standard monotone comparative statics
result---increasing differences imply increasing optimal
selections---adapted to the LR order on beliefs; it is a special case
of Lemma~6.1 of \citet{topkis1978minimizing}, and parallels Lemma~2.1
of \citet{lovejoy1987some}. \Claude{The statement of part (ii) below
was drafted by me---the notes contained only the placeholder
``increasing differences to increasing policies.'' Please verify the
statement (in particular, the use of the greatest maximizer to obtain a
well-defined monotone selection).}

\begin{lemma}\label{lemma:incdiff}\quad
\begin{enumerate}
\item If $u(\theta,a)$ has increasing differences, then
  $u(\pi,a) = \int_{\theta \in \Theta} u(\theta,a) \pi(\theta)
  \dif \theta$ has LR-increasing differences.
\item If $u(\pi,a)$ has LR-increasing differences, then the greatest
  maximizer $\alpha^u(\pi) = \max\, \mathrm{argmax}_{a \in A}\,
  u(\pi,a)$ is LR-increasing: $\pi_2 \LR \pi_1$ implies
  $\alpha^u(\pi_2) \geq \alpha^u(\pi_1)$. In particular, if
  $L_\iX(x|\theta,a)$ and $g(\theta_{k-1}|\theta_k,a)$ satisfy the MLR
  property, then for any prior $\pi$ and action $a$, the induced signal
  selection $\alpha^u(\Pi_\iX(\pi,x,a))$ is increasing in the observed
  signal $x$.
\end{enumerate}
\end{lemma}

The second claim in part (ii) follows by combining the first claim with
Proposition~\ref{prop:monotone_beliefs}(ii): more favorable signals
lead to LR-larger posteriors and hence to larger selected actions. We
can now state the monotone policy result. As with
Proposition~\ref{prop:monotone_values}, the proof of part (i) is by
induction on $k$---the reward term has LR-increasing differences by
Lemma~\ref{lemma:incdiff}(i) and the continuation term satisfies
\eqref{eq:stochincdiff} with $u = V_{\iX,k-1}^*$, which is
LR-increasing by Proposition~\ref{prop:monotone_values}---and part (ii)
follows from part (i) via Lemma~\ref{lemma:incdiff}(ii). We omit the
details. \Claude{Confirm you want this proof sketch in text rather
than a full proof in an appendix; OR will ultimately expect complete
proofs in an e-companion.}

\begin{proposition}[Monotonicity of optimal policies]\label{prop:monotone_policies}
Suppose (i) $r(\theta,a)$ is increasing in $\theta$ for all $a$ and has
increasing differences, and (ii) $L_\iX(x|\theta,a)$ and
$g(\theta_{k-1}|\theta_k,a)$ satisfy the MLR property for all $a$ and
satisfy the stochastic version of increasing differences. Then:
\begin{enumerate}
\item $V_{\iX,k}(\pi,a)$ has LR-increasing differences for all $k$.
\item There exists an optimal policy $\boldsymbol{\alpha}_{\iX}^*$ that
  is LR-increasing.
\end{enumerate}
\end{proposition}

Propositions \ref{prop:monotone_values} and
\ref{prop:monotone_policies} describe the class of models we will call
\emph{monotone POMDPs}: models in which values increase and optimal
actions shift upward as beliefs improve in the LR order. As discussed
in the introduction, many POMDP models in operations research fall in
this class. In the remainder of the paper, we study when one signal
process is more valuable than another in such models.
